\documentclass[11pt]{article}
\usepackage[a4paper,margin=27mm]{geometry}
\usepackage{amsmath,amssymb,amsthm,mathtools,bm,microtype}
\usepackage[colorlinks=true,allcolors=blue]{hyperref}
\usepackage{booktabs,enumitem,array}
\setlist{nosep,leftmargin=*,topsep=2pt}
\newtheorem{theorem}{Theorem}[section]
\newtheorem{proposition}[theorem]{Proposition}
\newtheorem{lemma}[theorem]{Lemma}
\newtheorem{corollary}[theorem]{Corollary}
\theoremstyle{remark}
\newtheorem{remark}[theorem]{Remark}
\newcommand{\Herm}{\operatorname{Herm}}
\newcommand{\Tr}{\operatorname{Tr}}
\newcommand{\id}{\mathbf 1}
\newcommand{\CP}{\mathrm{CP}}
\newcommand{\cL}{\mathcal L}
\title{Operational Lorentz Geometry from Resolved Quantum Instruments\\[0.4em]\large Normalization, coherent refinement, and spectral references}
\author{}
\date{29 October 2026}

\hypersetup{pdfauthor={},pdftitle={Operational Lorentz Geometry from Resolved Quantum Instruments}}
\begin{document}
\maketitle

\begin{abstract}
Resolved quantum instruments expose information that is lost when a selected
filter is treated as a deterministic geometric transport.
Starting from the established qubit--Lorentz cone correspondence, we give an
endpoint-order-unit normalization that is covariant under independent passive
spin-frame changes, identify a relative phase fibre invisible to diagonal
instrument data, and distinguish nonmultiplicative maximal postselection from
an exactly refinable complete dephasing tree.
A second identifiability obstruction concerns spectral references: an
assignment on positive matrices in all finite dimensions that respects unitary
conjugation and direct sums and forgets scalar spectral shifts must be
constant.
The familiar spectral-floor centering violates disconnected-component
consistency, while retaining an absolute generator or supplying a calibrated
reference admits an explicit positive absorbing completion.
All principal statements are finite-dimensional and have analytic proofs;
portable symbolic checks reproduce the decisive identities and a finite
endpoint-covariance example.
The paper supplies operational constructions and limits on reconstruction,
without deriving a spacetime coframe, gravitational action or continuum
quantum dynamics.
\end{abstract}

\section{Problem and scope}

Cartan gravity distinguishes a Lorentz connection from a solder form and, after a
$G/H$ reduction, packages both into a de Sitter, anti-de Sitter, or Poincar\'e connection
\cite{MacDowellMansouri,StelleWest,Wise}.  Quantum instruments distinguish a selected
operation from the complete record of mutually exclusive outcomes
\cite{DaviesLewis,ChiribellaNetworks}.  The two distinctions become entangled if a
non-trace-preserving qubit filter is used to represent a local Lorentz boost: the
noncompact transformation is visible only before state normalization, while the physical
instrument must also include the rejected outcome.

The qubit--Lorentz correspondence itself is standard
\cite{ArrighiPatricotCone,ArrighiPatricotLorentz,VerstraeteLorentz,FullwoodLorentz}.
Our question concerns normalization, coherent composition and missing references.
The generic Lorentz measurement lift is already known; our use of it retains the
resolved rejection cost and separates that fact from endpoint covariance.
The same earlier work observes that transformed effects need not sum to a fixed
identity~\cite{ArrighiPatricotLorentz}.
Our question is narrower and operational.  Which parts of a Lorentz/Cartan reading remain
well defined after probabilities, endpoint frames, coherent phases, and refinement are
included?  The answer is a mixture of constructions and obstructions.  Table
\ref{tab:claims} states the evidence boundary.

\begin{table}[ht]
\centering
\begin{tabular}{@{}>{\raggedright\arraybackslash}p{.24\linewidth}>{\raggedright\arraybackslash}p{.34\linewidth}>{\raggedright\arraybackslash}p{.32\linewidth}@{}}
\toprule
Statement & Status & Not established \\
\midrule
$\Herm(2)_+$ is the future cone & background theorem & spacetime soldering \\
resolved filter lift & analytic construction & gravitational transport law \\
moving-order-unit covariance & analytic theorem and finite check & microscopic clock selection \\
Kraus-phase blindness & analytic no-go & superselection of the phase \\
complete-tree refinement & exact special channel & cellulation independence \\
\bottomrule
\end{tabular}
\caption{Claim boundary.  ``Background'' is not an originality claim.}
\label{tab:claims}
\end{table}

We work in finite dimensions.  Effects and unnormalised states are positive Hermitian
matrices.  The Schr\"odinger branch is $\rho\mapsto K\rho K^\dagger$; with a local order
unit $u>0$, its Heisenberg normalization is $K^\dagger u_{\rm out}K\leq u_{\rm in}$.
No continuum limit or causal signaling theorem is assumed.

\section{The qubit order cone and the Lorentz metric}

Let $J=\Herm(2)$ as a four-dimensional real vector space.  In a Pauli frame,
\begin{equation}
 X=x^0\id+x^i\sigma_i,
 \qquad q(X)=\det X=(x^0)^2-\bm x^2 .
 \label{eq:minkowski-determinant}
\end{equation}
Polarization gives
\begin{equation}
 \eta(X,Y)=\frac12\{q(X+Y)-q(X)-q(Y)\}
 =\frac12\bigl(\Tr X\Tr Y-\Tr(XY)\bigr).
 \label{eq:det-polarization}
\end{equation}

\begin{proposition}[Order cone]
\label{prop:order-cone}
The positive cone is
\begin{equation}
 J_+=\{X\geq0\}=\{X:q(X)\geq0,\ \Tr X\geq0\}.
 \label{eq:future-cone}
\end{equation}
Rank-one positive matrices are future null and positive-definite matrices are future
timelike.
\end{proposition}

\begin{proof}
The eigenvalues are $x^0\pm|\bm x|$.  Simultaneous nonnegativity is therefore
$x^0\geq|\bm x|$, while their product is \eqref{eq:minkowski-determinant}.
\end{proof}

For $S\in SL(2,\mathbb C)$, $X\mapsto SXS^\dagger$ preserves the cone and determinant.
Modulo $\{\pm\id\}$ this is the standard spin cover of $SO^+(1,3)$.  The singlet
correlation also reads out \eqref{eq:det-polarization}:
\begin{equation}
 \langle\psi_-|X\otimes Y|\psi_-\rangle
 =x^0y^0-\bm x\cdot\bm y=\eta(X,Y).
 \label{eq:singlet-metric}
\end{equation}
This last observation overlaps the explicit analysis of Ref.~\cite{FullwoodLorentz}; we
use it only as a basis-independent readout.  It does not identify $J$ with a tangent
space.

There is a useful negative statement.  If $\rho:M\to\mathcal S_2^+$ is a faithful
trace-one state field and $G$ is any positive monotone quantum-information metric, then
$G_\rho(d\rho(v),d\rho(v))\geq0$.  No real change of coordinates turns its nondegenerate
pullback into a Lorentz metric.  This follows directly from positivity and Sylvester
inertia and applies, for example, to the Petz family \cite{PetzMonotone}.  The indefinite
form in \eqref{eq:det-polarization} lives on unnormalised observables, not on the tangent
space of normalized states.

\subsection{Representability is not derivation}

Any four Hermitian matrices $\Theta_\mu$ can be encoded in four independently selected
binary settings,
\begin{equation}
 E_{\mu,\pm}(a)=\frac12\id\pm a\Theta_\mu,
 \qquad
 \Theta_\mu=\frac{E_{\mu,+}-E_{\mu,-}}{2a},
 \label{eq:balanced-effects}
\end{equation}
provided $|a|\|\Theta_\mu\|_{\rm op}\leq1/2$.  If a microscopic law independently
supplies a four-dimensional deformation space $W_x$ and a linear full-rank response
$\Theta_x:W_x\to J$, then
\begin{equation}
 g_x(v,w)=\eta(\Theta_xv,\Theta_xw)
 \label{eq:solder-pullback}
\end{equation}
is Lorentzian.  Equation \eqref{eq:balanced-effects} proves encodability only.  Because an
arbitrary coframe can be inserted, it cannot explain its rank, scale, anholonomy, or
dynamics.

\begin{lemma}[Operational future cone]
\label{lem:activation}
For $X\in J$, $X\geq0$ if and only if there is an $\epsilon>0$ such that
$\{aX,\id-aX\}$ is a binary POVM for every $0\leq a\leq\epsilon$.
\end{lemma}

\begin{proof}
If the pair is positive for any $a>0$, then $X\geq0$.  Conversely choose
$\epsilon\leq\|X\|_{\rm op}^{-1}$ for positive nonzero $X$.
\end{proof}

Thus, conditional on a valid solder response, the future cone consists exactly of
directions that can activate a previously absent outcome to first order.  This is a local
order-theoretic orientation, not yet a domain-of-dependence statement.

\section{Resolved filters and the cost of boosts}

The scaled measurement realization below is established in Proposition~3 of Ref.~\cite{ArrighiPatricotLorentz}; we include the proof to fix its completion and normalization conventions.

\begin{theorem}[Filter factorization and completion]
\label{thm:filter-lift}
Let $G\in SL(2,\mathbb C)$.  For every
$0<r\leq\|G\|_{\rm op}^{-1}$, $K=rG$ is a trace-nonincreasing single-Kraus branch and
\begin{equation}
 L=(\id-K^\dagger K)^{1/2}
 \label{eq:complement-kraus}
\end{equation}
completes it to a two-outcome trace-preserving instrument.  Conversely, every invertible
$K$ factorizes locally as $K=cG$, with $G\in SL(2,\mathbb C)$ and
$c=(\det K)^{1/2}$, up to the spin-cover sign.
\end{theorem}

\begin{proof}
The norm bound gives $K^\dagger K\leq\id$, so \eqref{eq:complement-kraus} is positive and
$K^\dagger K+L^\dagger L=\id$.  The converse follows from
$G=c^{-1}K$ and $\det G=1$ on any continuous square-root branch.
\end{proof}

The scalar $r^2$ is not the branch probability.  For a normalized input,
\begin{equation}
 p_K(\rho)=r^2\Tr(G^\dagger G\rho),
 \qquad q(KXK^\dagger)=r^4q(X).
 \label{eq:filter-probability}
\end{equation}
The input dependence is precisely why discarding the rejected outcome erases operational
information.

Let the singular values of $G$ be $s,s^{-1}$ with $s\geq1$.  The two contractions
\begin{equation}
 K_+=s^{-1}G,\qquad K_-=s^{-1}G^{-1}
 \label{eq:two-way-filters}
\end{equation}
obey $K_-K_+=K_+K_-=s^{-2}\id$.  Projectively the selected branches are inverse Lorentz
maps, but their unnormalised round trip is the central loss $X\mapsto s^{-4}X$.

\begin{corollary}[Null rejected outcome for a pure boost]
\label{cor:null-failure}
For
$G_\chi=\exp[(\chi/2)\bm n\cdot\bm\sigma]$, $\chi>0$, the maximally scaled branch
$K=e^{-\chi/2}G_\chi$ has
\begin{equation}
 \operatorname{spec}(K^\dagger K)=\{1,e^{-2\chi}\},\qquad
 \id-K^\dagger K=(1-e^{-2\chi})P_-,
 \label{eq:null-failure}
\end{equation}
where $P_-$ is rank one.  The rejected effect is therefore a future-null element of
$J_+$.
\end{corollary}

The result relates rapidity to an exact survival cost, but it does not identify boost loss
with gravitational redshift.  That interpretation requires endpoint clocks and a
refinement law.

\section{Moving order units and passive covariance}

Using a fixed $\id$ as the order unit at every endpoint confuses a passive nonunitary
change of spin frame with a physical filter.  Starting with an identity edge and applying
different $S_x,S_y\in SL(2,\mathbb C)$ gives $G'=S_yS_x^{-1}$; normalizing by the fixed
operator norm generally creates a nonzero complement.  Outcome probabilities then depend
on a gauge choice.

Assign instead an independently reconstructed $u_x>0$ to each endpoint.  Effects satisfy
$0\leq E\leq u_x$, and dual states obey $\Tr(\tau_xu_x)=1$.  Under a passive frame change,
\begin{equation}
 E\mapsto S_xES_x^\dagger,\quad
 u_x\mapsto S_xu_xS_x^\dagger,\quad
 \tau_x\mapsto S_x^{-\dagger}\tau_xS_x^{-1}.
 \label{eq:order-unit-frame}
\end{equation}
Both the order interval and probability pairing are invariant.

\begin{theorem}[Order-unit covariant resolved instrument]
\label{thm:order-unit-instrument}
For an invertible edge matrix $G_{xy}$ and positive endpoint units, define
\begin{align}
 A_{xy}&=u_x^{-1/2}G_{xy}^{-1}u_yG_{xy}^{-\dagger}u_x^{-1/2},\nonumber\\
 c_{xy}&=\lambda_{\max}(A_{xy})^{-1/2},\qquad
 K_{g,xy}=c_{xy}G_{xy}^{-\dagger}.
 \label{eq:covariant-normalization}
\end{align}
Then
\begin{equation}
 K_{g,xy}^\dagger u_yK_{g,xy}\leq u_x,
 \qquad E_n=u_x-K_{g,xy}^\dagger u_yK_{g,xy}\geq0.
 \label{eq:order-unit-complement}
\end{equation}
Under independent endpoint changes
$G'_{xy}=S_yG_{xy}S_x^{-1}$ and $u'_z=S_zu_zS_z^\dagger$, the selected and rejected
effects transform by source congruence and all paired probabilities are invariant.
\end{theorem}

\begin{proof}
The largest-eigenvalue choice makes $c_{xy}^2A_{xy}\leq\id$, which is congruent to the
first inequality in \eqref{eq:order-unit-complement}.  Generalized eigenvalues are
unchanged by the simultaneous endpoint congruences.  Direct substitution gives
$K'_g=S_y^{-\dagger}K_gS_x^\dagger$, so
$K_g'^\dagger u_y'K_g'=S_x(K_g^\dagger u_yK_g)S_x^\dagger$ and the same holds for
$E_n$.
\end{proof}

In standardized fibers, $L_g=u_y^{1/2}K_gu_x^{-1/2}$ transforms as
$L'_g=R_yL_gR_x^\dagger$ for unitary Wigner-frame matrices $R_x,R_y$.  Its singular
values are therefore frame-independent mismatch data.  A local Kraus completion always
exists, for example $K_n=u_y^{-1/2}E_n^{1/2}$, but the effect $E_n$ does not determine the
output partial isometry or coherent phase of $K_n$.

A deterministic special case deserves care.  If an independent measurement verifies
\begin{equation}
 u_y=G_{xy}u_xG_{xy}^\dagger,
 \label{eq:moving-unit-lock}
\end{equation}
then $U_{xy}=u_y^{-1/2}G_{xy}u_x^{1/2}$ is unitary.  Noncompact information is stored in
the change of endpoint unit, while standardized transport is compact.  Defining $u_y$ by
\eqref{eq:moving-unit-lock} after measuring $G$ would make the statement tautological.
Moreover, imposing it around every loop forces holonomy into the compact stabilizer of
$u_x$, excluding generic Lorentz curvature.  It is therefore a calibrated subfamily, not
a universal gravitational law.

\section{Coherent phase is not instrument data}

Let a labeled instrument have one nonzero Kraus operator in each outcome,
\begin{equation}
 \mathcal I_a(\rho)=K_a\rho K_a^\dagger,\qquad
 \sum_aK_a^\dagger K_a=\id.
 \label{eq:labeled-instrument}
\end{equation}

\begin{theorem}[Diagonal phase blindness]
\label{thm:phase-blindness}
The transformations $K_a\mapsto e^{i\theta_a}K_a$ leave every labeled CP branch, branch
probability, conditioned output, and observable diagonal in the record basis invariant.
After a common phase is removed, an $m$-outcome single-Kraus instrument has an
unidentified $U(1)^{m-1}$ phase fiber whenever all outcomes remain labeled.
\end{theorem}

\begin{proof}
The phase cancels in $K_a\rho K_a^\dagger$.  For the Stinespring isometry
$V=\sum_a|a\rangle\otimes K_a$ and $D=\sum_ad_a|a\rangle\langle a|$,
\begin{equation}
 V^\dagger(D\otimes O)V=\sum_ad_aK_a^\dagger OK_a,
 \label{eq:diagonal-record}
\end{equation}
which contains no cross term.  A simultaneous rephasing of all $K_a$ is global.
\end{proof}

The phase is not superselected.  An off-diagonal record observable gives
\begin{equation}
 V^\dagger(|a\rangle\langle b|\otimes O)V=K_a^\dagger OK_b,
 \label{eq:coherent-record}
\end{equation}
which acquires $e^{i(\theta_b-\theta_a)}$.  A coherent recombiner can therefore measure a
relative phase if the matrix element is nonzero.  The theorem says that diagonal
instrument tomography cannot select which coherent dilation nature uses.

This distinction matters for gravity: classical equations are invariant under an overall
rescaling $S\mapsto cS$ of the action, while $\exp(iS/\hbar)$ is not.  Even a finite record
that verifies Einstein--Cartan stationary equations would not by itself determine Newton's
constant or the phase accumulated between alternative geometries.

\section{Refinement: selected branches versus complete trees}
\label{sec:refinement}

Suppose a Hermitian mismatch $X$ determines the locally maximized selected contraction
rate
\begin{equation}
 R=\lambda_{\max}(X)\id-X\geq0,\qquad
 K_0(dt)=\id-\tfrac12R\,dt+O(dt^2).
 \label{eq:selected-rate}
\end{equation}
If each subedge is independently maximized, the normalization depends nonlinearly on the
subedge data.  In general the product of fine selected branches is not the maximized coarse
branch.  This is a functoriality obstruction for postselection, not a failure of CP
composition.

The correct normalized object is the complete process.  Adding $J=\sqrt R$ gives
\begin{equation}
 \cL(\rho)=J\rho J^\dagger-	frac12\{J^\dagger J,\rho\}
 \label{eq:gksl-generator}
\end{equation}
up to an optional Hamiltonian commutator.  For the rank-one pure-boost mismatch
$R=2bP_+$ with $b\geq0$, an exact two-Kraus channel is
\begin{equation}
 K_0(t)=e^{-bt}P_++P_-,\qquad
 K_1(t)=\sqrt{1-e^{-2bt}}\,P_+.
 \label{eq:exact-channel}
\end{equation}

\begin{theorem}[Exact complete-tree refinement]
\label{thm:tree-refinement}
For $t_1,t_2\geq0$, the four fine histories generated by
\eqref{eq:exact-channel} can be unitarily recombined in record space into the two coarse
Kraus operators at $t_1+t_2$.  Repeating the construction with positive real time-bin
coefficients gives an associative record isometry for any finite ordered subdivision.
\end{theorem}

\begin{proof}
Write $s_i=e^{-bt_i}$ and $j_i=\sqrt{1-s_i^2}$.  The four histories are
\begin{align}
 K_{00}&=s_1s_2P_++P_-,&K_{01}&=s_2j_1P_+,\nonumber\\
 K_{10}&=s_1j_2P_+,&K_{11}&=j_1j_2P_+.
 \label{eq:fine-histories}
\end{align}
The three jump histories are collinear, and their squared coefficients sum to
\begin{equation}
 s_2^2j_1^2+s_1^2j_2^2+j_1^2j_2^2=1-s_1^2s_2^2.
 \label{eq:jump-norm}
\end{equation}
A unitary record change therefore compresses them to $K_1(t_1+t_2)$, while $K_{00}$ is
already $K_0(t_1+t_2)$.  For more intervals, the vector of nonzero jump-history
coefficients has norm $\sqrt{1-e^{-2b\sum_it_i}}$.  Normalizing this vector defines the
record isometry independently of parenthesization, proving associativity in the declared
positive time-bin gauge whenever the total jump weight is nonzero.
If all jump weights vanish, only the deterministic no-jump record is supported; its
embedding is the all-zero history and the unused zero-amplitude record can be omitted.
\end{proof}

The channel does not select that gauge.  Rephasing one fine jump history leaves the CP map
unchanged but changes a fixed coherent recombiner, exactly as Theorem
\ref{thm:phase-blindness} predicts.  Nor does one-dimensional chain refinement imply
equivalence of competing spacetime cellulations, constraint closure, or a continuum
amplitude.

\section{Spectral references and disconnected components}
\label{sec:spectral}

Phase is not the only datum that normalized operational records can forget.
The maximal mismatch rule of Section~\ref{sec:refinement} assigns a positive rate
$R=\lambda_{\max}(X)I-X$ to a Hermitian mismatch $X$.
Given a positive operator $H$, the choice $X=\lambda_{\max}(H)I-H$ yields
\begin{equation}
 R=C(H):=H-\lambda_{\min}(H)I.
 \label{eq:centering}
\end{equation}
This is invariant under every admissible shift $H\mapsto H+cI$.
Here $H$ has rate units, or the readout parameter has been calibrated to supply
the reciprocal units; the notation alone supplies no physical clock.
The no-jump amplitude $e^{-tR/2}$ and its survival probabilities therefore cannot
identify the absolute spectral floor.

\begin{proposition}[Disconnected-spectator obstruction]
\label{prop:centering}
For positive matrices $H_1,H_2$ of nonzero dimensions, let
$\lambda_i=\lambda_{\min}(H_i)$.
Then
\begin{equation}
 C(H_1\oplus H_2)|_i
 =C(H_i)+\bigl(\lambda_i-\min\{\lambda_1,\lambda_2\}\bigr)I_i.
 \label{eq:centering-defect}
\end{equation}
Consequently $C$ preserves this direct sum if and only if
$\lambda_1=\lambda_2$.
\end{proposition}
\begin{proof}
The spectrum of a direct sum is the union of its component spectra, so its
floor is the minimum of their floors.
Subtracting that common floor from each block gives
\eqref{eq:centering-defect}; both offsets vanish precisely in the stated case.
\end{proof}

For example, take
\begin{equation}
 H_1=\operatorname{diag}(1/5,11/10),\qquad
 H_2=\operatorname{diag}(7/10,3/2).
\end{equation}
The two results are
\begin{align}
 C(H_1)\oplus C(H_2)&=\operatorname{diag}(0,9/10,0,4/5),\nonumber\\
 C(H_1\oplus H_2)&=\operatorname{diag}(0,9/10,1/2,13/10).
\end{align}
The disconnected lower-floor component adds $I/2$ to the rate on the second
component.
At a fixed readout time, its survival probability is multiplied by $e^{-t/2}$.
Thus this defect is observable in a resolved no-jump record; it is not merely
an overall Hamiltonian phase.
Direct-sum preservation is the precise disconnected-component requirement here.
It should not be equated with every notion of relativistic locality or with
tensor-product composition.

\subsection{A general finite reference theorem}

Changing the scalar statistic in \eqref{eq:centering} cannot preserve all of the
same requirements while retaining spectral information.
The following statement makes this obstruction independent of the choice of
the minimum eigenvalue.

\begin{theorem}[Spectral-reference obstruction]
\label{thm:reference}
For every $n\geq1$, let
$F_n:\Herm(n)_+\to\Herm(n)$.
Assume unitary covariance, preservation of all finite direct sums, and scalar
shift invariance:
\begin{align}
 F_n(UHU^\dagger)&=UF_n(H)U^\dagger,
 \label{eq:reference-unitary}\\
 F_{n+m}(H\oplus K)&=F_n(H)\oplus F_m(K),
 \label{eq:reference-sum}\\
 F_n(H+cI)&=F_n(H)
 \label{eq:reference-shift}
\end{align}
whenever the arguments are positive semidefinite.
There is a real constant $a$, independent of $n,H$, such that
$F_n(H)=aI_n$.
If $F_1(0)=0$, every $F_n$ is zero.
No continuity, linearity, differentiability or spectral nondegeneracy
assumption is required.
\end{theorem}
\begin{proof}
Write $F_1([\lambda])=[f(\lambda)]$.
Equation~\eqref{eq:reference-shift} applied to $[0]$ and $[\lambda]$ gives
$f(\lambda)=f(0)=a$ for every $\lambda\geq0$.
By repeated use of \eqref{eq:reference-sum},
\[
 F_n(\operatorname{diag}(\lambda_1,\ldots,\lambda_n))
 =\operatorname{diag}(f(\lambda_1),\ldots,f(\lambda_n))=aI_n.
\]
The spectral theorem and \eqref{eq:reference-unitary} extend this equality to
every positive Hermitian $H$.
The final assertion follows from $a=F_1(0)$.
\end{proof}

The theorem is deliberately a statement about bare matrices in all finite
dimensions.
A model equipped with a preferred component algebra has additional data and
need not be covariant under every unitary that mixes those components.
Componentwise centering can then be defined relative to that algebra, but the
component choice has not been obtained from a shift-blind matrix record.
Replacing the minimum by a mean or quantile does not evade
Theorem~\ref{thm:reference} under its hypotheses.

\subsection{Calibrated repairs and a complete positive process}

There are two elementary repairs.
One may retain $H$ itself and its absolute spectral floor, thereby preserving
unitary covariance and direct sums while abandoning shift blindness.
Alternatively supply a reference operator $B$ with $R=H-B\geq0$.
The assignment
\begin{equation}
 R(H,B)=H-B,\qquad
 R(H+cI,B+cI)=R(H,B)
 \label{eq:calibrated-reference}
\end{equation}
preserves direct sums of pairs and simultaneous unitary transformations.
It contains an extra calibrated datum; it does not extract an absolute
generator from the original normalized record.

A complete process can be attached to any chosen $R\geq0$ without imposing a
two-Kraus chain law on a general spectrum.
Let $R=\sum_{j=1}^d r_j|j\rangle\langle j|$, introduce an absorbing flag
$|\varnothing\rangle$ orthogonal to the live space, and define
\begin{align}
 A_0(t)&=|\varnothing\rangle\langle\varnothing|
       +\sum_j e^{-tr_j/2}|j\rangle\langle j|,\nonumber\\
 A_j(t)&=\sqrt{1-e^{-tr_j}}\,
           |\varnothing\rangle\langle j|,\qquad j=1,\ldots,d .
 \label{eq:absorbing-completion}
\end{align}

\begin{proposition}[Calibrated absorbing completion]
\label{prop:absorbing}
The operators \eqref{eq:absorbing-completion} define a trace-preserving
completely positive semigroup on the live space plus the absorbing flag.
For a normalized live input $\rho$, its probability of remaining live is
$\Tr(e^{-tR}\rho)$.
For the maximally mixed live input this is $d^{-1}\Tr e^{-tR}$.
\end{proposition}
\begin{proof}
The completeness identity follows from
$e^{-tr_j}+(1-e^{-tr_j})=1$ for each live mode and the unit coefficient on the
absorbing flag.
Live matrix elements evolve by
$\rho_{jk}\mapsto e^{-t(r_j+r_k)/2}\rho_{jk}$, live--flag coherences by
$\rho_{j\varnothing}\mapsto e^{-tr_j/2}\rho_{j\varnothing}$, and the flag
population receives $\sum_j(1-e^{-tr_j})\rho_{jj}$.
These formulas compose at times $t_1,t_2$ by addition of the times.
Equivalently, the generator has jump operators
$L_j=\sqrt{r_j}|\varnothing\rangle\langle j|$.
Taking the trace of the live block gives the asserted probabilities.
\end{proof}

This construction is a calibrated measurement of a heat trace.
For a geometric choice $H=D^\dagger D$, convergence to a continuum spectrum
and the physical meaning of the readout time remain separate questions.
A protected zero mode anchors the floor only in sectors that contain it.
It is not universal: on the round $S^4_r$ the untwisted self-adjoint Dirac
operator has eigenvalues $\pm(k+2)/r$, so $D^2$ has floor $4/r^2$
\cite{Camporesi}.
The floor belongs to $D^2$, not to the two-sided Dirac spectrum itself.
Neither an absorbing completion nor a chosen reference $B$ supplies a
gravitational action or selects a microscopic clock.


\section{Cartan-geometric interpretation and falsifiers}

The results support a disciplined operational dictionary:
\begin{enumerate}
\item $\Herm(2)_+$ supplies an internal causal cone and Lorentz bilinear form.
\item A separately derived full-rank response is required to solder this
internal space to spacetime; curved geometries generally require anholonomic coframes.
\item Resolved filters can tomographically infer noncompact transport, but passive frame
covariance requires moving endpoint order units.
\item Rejected outcomes carry noncompact survival data; they cannot be silently removed.
\item Quantum composition requires a coherent record structure beyond diagonal CP data.
\end{enumerate}

Each item has a direct falsifier.  Failure of response linearity or rank kills the coframe
reading.  Frame-dependent generalized singular values kill passive covariance.  A coarse
instrument not obtainable from the complete fine tree kills the proposed refinement.  An
interference experiment inconsistent with the declared record isometry kills its coherent
lift.  None of these tests presupposes the Einstein equations.

Several tempting promotions are excluded.  Gradients of one four-component coordinate
field, $e^I=du^I$, give a flat metric when full rank and a degenerate metric otherwise; a
curved coframe must be generally anholonomic.  Derivatives of normalized qubit probabilities
span at most the three traceless directions, so the trace/time component requires an
independently registered rate or scale.  Finally, a finite-dimensional unitary
representation of the connected noncompact simple group $SL(2,\mathbb C)$ is trivial at
the Lie-algebra level; adding a finite record register cannot turn the whole Lorentz gauge
group into unitary transformations inside one fixed fiber.

\section{Conclusion}

The standard qubit Lorentz cone becomes operationally consistent only when the full
instrument, its endpoint order units, and its coherent record are kept distinct.  The
selected filter carries the projective Lorentz map; the rejected branch carries its
survival cost; moving order units separate passive frame changes from physical loss; and a
declared coherent record isometry, not diagonal channel data, carries refinement phases.
This supplies finite Lorentz--Cartan kinematics and no-go theorems about identifiability.
Theorem~\ref{thm:reference} adds a separate boundary on spectral reconstruction: a
nontrivial direct-sum-covariant record cannot forget every absolute scalar reference.
The calibrated completion in Proposition~\ref{prop:absorbing} is an existence construction,
not a deduction of a physical generator.
The missing steps are substantive: a microscopic full-rank coframe, a dynamical
Lorentz/Cartan connection, a gravitational amplitude with constraints, and a controlled
continuum limit.  We therefore make no claim of emergent Einstein dynamics or quantum
gravity.

\appendix
\section{Portable verification and provenance}
The companion program \texttt{code/verify.py} runs without any historical project
path or network access.
SymPy checks the boost complement, the opposite-boost composition defect, the
jump-history polynomial identity, a coherent recombiner with probabilities one
and zero for identical diagonal instruments, and the exact rational
spectral-floor witness.
A finite NumPy example with noncommuting endpoint data checks the scale,
selected and rejected effects, and probability covariance; its largest
residual is below $3\times10^{-16}$.
Small negative rejected-effect eigenvalues of order $10^{-16}$ in this example
are roundoff, not an analytic positivity proof.
The proofs in the body establish the general claims.
No continuum, gravitational field-equation, or all-input numerical-stability
claim is inferred from these finite checks.

The underlying local research records are dated 12--13 July 2026; this
English edition and the spectral-reference strengthening were prepared on
29 October 2026.
Private source dates do not establish public disclosure or priority.
The operator reports that earlier work was led by Fable~5, while results from
July 2026 onward were predominantly GPT-produced; DeepSeek and multiple GPT
models also contributed review and preparation.
Human authors and affiliations are left blank.
Human expert review has not occurred, and no global novelty certificate is
claimed.
The operator's attribution statement is retained verbatim:
\begin{quote}
I genuinely do not know how authorship, intellectual credit, or ownership
should be attributed for these results.
\end{quote}

\bibliographystyle{unsrt}
\bibliography{references}
\end{document}
